\chapter{1988年广东卷（理）}

\section{单选题}

\begin{problem}
直线 \(3x - 2y = 6\) 在 \(y\) 轴上的截距是（\quad）

\choices
    {\(\frac{3}{2}\)}
    {\(-2\)}
    {\(-3\)}
    {\(3\)}
\end{problem}

\begin{answer}
C
\end{answer}

\begin{solution}
在 \(y\) 轴上有 \(x=0\)，代入直线方程得 \(-2y=6\)，所以 \(y=-3\). 因此直线在 \(y\) 轴上的截距为 \(-3\)，故选 C.
\end{solution}

\begin{problem}
对数方程 \(2\log_6 x = 1 - \log_6 3\) 的解是（\quad）

\choices
    {\(\sqrt{3}\)}
    {\(\sqrt{2}\)}
    {\(\sqrt{2}\) 和 \(-\sqrt{2}\)}
    {\(\sqrt{3}\) 和 \(\sqrt{2}\)}
\end{problem}

\begin{answer}
B
\end{answer}

\begin{solution}
由定义域知 \(x>0\). 原方程可化为 \(2\log_6x=\log_6 6-\log_6 3=\log_6 2\)，即 \(\log_6x^2=\log_6 2\). 因而 \(x^2=2\)，结合 \(x>0\) 得 \(x=\sqrt{2}\)，故选 B.
\end{solution}

\begin{problem}
如果用 \(C\)，\(\R\) 和 \(I\) 分别表示复数集，实数集和纯虚数集，其中 \(C\) 为全集，那么有（\quad）

\choices
    {\(C = \R \cup I\)}
    {\(\R \cap I = \{0\}\)}
    {\(\overline{\R} = C \cap I\)}
    {\(\R \cap I = \varnothing\)}
\end{problem}

\begin{answer}
D
\end{answer}

\begin{solution}
按题设，纯虚数集 \(I\) 不含 \(0\)，而实数集 \(\R\) 中的数都不是纯虚数，所以 \(\R\cap I=\varnothing\)，选项 B 错误. 取 \(1+\i\in C\)，它既不是实数也不是纯虚数，所以 \(C\ne\R\cup I\)，选项 A 错误；同时 \(1+\i\) 是非实复数，故 \(1+\i\in\overline{\R}\)，但 \(1+\i\notin I\)，所以 \(\overline{\R}=C\cap I\) 也不成立. 因此选 D.
\end{solution}

\begin{problem}
在等差数列 \(\{a_n\}\) 中，已知前 \(15\) 项之和 \(S_{15} = 90\)，那么 \(a_8 =\)（\quad）

\choices
    {\(3\)}
    {\(4\)}
    {\(6\)}
    {\(12\)}
\end{problem}

\begin{answer}
C
\end{answer}

\begin{solution}
等差数列前 \(15\) 项关于第 \(8\) 项成对对称，因此 \(S_{15}=15a_8\). 由 \(15a_8=90\) 得 \(a_8=6\)，故选 C.
\end{solution}

\begin{problem}
如果 \(\tan A + \cot A = m\)，那么 \(\sin 2A =\)（\quad）

\choices
    {\(\frac{1}{m}\)}
    {\(\frac{2}{m}\)}
    {\(2m\)}
    {\(\frac{1}{m^2}\)}
\end{problem}

\begin{answer}
B
\end{answer}

\begin{solution}
因为 \(\tan A\) 与 \(\cot A\) 有定义，所以 \(\sin A\cos A\ne0\). 因而
\(m=\tan A+\cot A=\dfrac{\sin^2A+\cos^2A}{\sin A\cos A}=\dfrac{2}{\sin2A}\)，从而 \(\sin2A=\dfrac{2}{m}\)，故选 B.
\end{solution}

\begin{problem}
设 \(a\)，\(b\) 是异面直线，那么（\quad）

\choices
    {必然存在唯一的一个平面同时平行于直线 \(a\) 和 \(b\)}
    {必然存在唯一的一个平面同时垂直于直线 \(a\) 和 \(b\)}
    {过直线 \(a\) 存在唯一的一个平面平行于直线 \(b\)}
    {过直线 \(a\) 存在唯一的一个平面垂直于直线 \(b\)}
\end{problem}

\begin{answer}
C
\end{answer}

\begin{solution}
取 \(A\in a\)，过 \(A\) 作唯一的直线 \(a'\parallel b\). 因为 \(a\)，\(b\) 是异面直线，\(a\) 与 \(a'\) 不平行，故它们确定唯一平面 \(\alpha\). 由 \(a'\subset\alpha\) 且 \(a'\parallel b\)，直线 \(b\) 的方向平行于平面 \(\alpha\). 若 \(b\) 与 \(\alpha\) 相交，则 \(b\) 必在 \(\alpha\) 内，这将使 \(a\)，\(b\) 共面，与异面矛盾. 所以 \(b\parallel\alpha\)，即存在这样的平面.

若另有平面 \(\beta\) 经过 \(a\) 且平行于 \(b\)，则过 \(A\) 作的平行于 \(b\) 的直线必须同时位于 \(\beta\) 内，因而就是唯一的 \(a'\). 所以 \(\beta\) 与 \(a\)，\(a'\) 共面，只能与 \(\alpha\) 重合. 因此该平面唯一，选 C.
\end{solution}

\begin{problem}
当 \(a \neq 0\) 时，方程 \(x^2 + y^2 + ax - ay = 0\) 所表示的图形是（\quad）

\choices
    {关于 \(x\) 轴对称}
    {关于 \(y\) 轴对称}
    {关于直线 \(y - x = 0\) 对称}
    {关于直线 \(y + x = 0\) 对称}
\end{problem}

\begin{answer}
D
\end{answer}

\begin{solution}
配方得 \(x^2+ax+y^2-ay=\left(x+\dfrac a2\right)^2+\left(y-\dfrac a2\right)^2-\dfrac{a^2}{2}=0\). 这是圆心为 \(\left(-\dfrac a2,\dfrac a2\right)\) 的圆，且圆心在直线 \(x+y=0\) 上；关于该直线反射时圆心不变，故图形关于 \(y+x=0\) 对称，选 D.
\end{solution}

\begin{problem}
设 \(\theta\) 是第四象限的角，那么方程 \((\sin \theta)x^2 + y^2 = \sin 2\theta\) 所表示的曲线是（\quad）

\choices
    {焦点在 \(x\) 轴上的椭圆}
    {焦点在 \(y\) 轴上的椭圆}
    {焦点在 \(x\) 轴上的双曲线}
    {焦点在 \(y\) 轴上的双曲线}
\end{problem}

\begin{answer}
C
\end{answer}

\begin{solution}
第四象限内 \(\sin\theta<0\) 且 \(\cos\theta>0\)，所以 \(\sin2\theta=2\sin\theta\cos\theta<0\). 原方程化为
\(\left(-\sin\theta\right)x^2-y^2=-\sin2\theta\)，其中两边的系数均为正，故表示焦点在 \(x\) 轴上的双曲线，选 C.
\end{solution}

\begin{problem}
抛物线 \(y^2 = 4x - 2\) 的准线的方程是（\quad）

\choices
    {\(x = -\frac{1}{2}\)}
    {\(x = \frac{1}{4}\)}
    {\(x = 0\)}
    {\(x = -1\)}
\end{problem}

\begin{answer}
A
\end{answer}

\begin{solution}
将方程写成 \(y^2=4\left(x-\dfrac12\right)\). 与标准式 \(y^2=4p(x-h)\) 比较得 \(h=\dfrac12\)，\(p=1\)，所以准线为 \(x=h-p=-\dfrac12\)，故选 A.
\end{solution}

\begin{problem}
过抛物线 \(y^2 = 4x\) 的焦点，作直线与此抛物线相交于两点 \(P\) 和 \(Q\)，那么线段 \(PQ\) 中点的轨迹的方程是（\quad）

\choices
    {\(y^2 = 2x - 1\)}
    {\(y^2 = 2x - 2\)}
    {\(y^2 = -2x + 1\)}
    {\(y^2 = -2x + 2\)}
\end{problem}

\begin{answer}
B
\end{answer}

\begin{solution}
过焦点 \((1,0)\) 的直线若为 \(y=0\)，只与抛物线相切于顶点，不能得到两个交点；因此先设所作非竖直割线为 \(y=m(x-1)\)，其中 \(m\ne0\). 与 \(y^2=4x\) 联立，交点横坐标 \(x_1\)，\(x_2\) 满足

\(m^2(x-1)^2=4x\)，\(\Longrightarrow\)，\(m^2x^2-(2m^2+4)x+m^2=0\).

其判别式为 \(\Delta=(2m^2+4)^2-4m^4=16(m^2+1)>0\)，所以每个 \(m\ne0\) 都对应两个不同的交点. 由韦达定理，\(x_1+x_2=2+\dfrac4{m^2}\). 因而交点中点 \((X,Y)\) 满足 \(X=\frac{x_1+x_2}{2}=1+\frac2{m^2}\)，\(Y=\frac{m(x_1-1)+m(x_2-1)}2=\frac2m\).

消去 \(m\) 得 \(Y^2=\dfrac4{m^2}=2X-2\). 还需考虑竖直线 \(x=1\)：它与抛物线交于 \((1,2)\)、\((1,-2)\)，中点为 \((1,0)\)，也满足 \(Y^2=2X-2\). 因此轨迹方程为 \(y^2=2x-2\)，选 B.
\end{solution}

\begin{problem}
设正方体的全面积为 \(24\mathrm{~cm}^2\)，一个球内切于该正方体，那么，这个球的体积是（\quad）

\choices
    {\(\frac{4}{3}\pi\mathrm{~cm}^3\)}
    {\(\frac{8}{3}\pi\mathrm{~cm}^3\)}
    {\(\sqrt{6}\pi\mathrm{~cm}^3\)}
    {\(\frac{32}{3}\pi\mathrm{~cm}^3\)}
\end{problem}

\begin{answer}
A
\end{answer}

\begin{solution}
设正方体棱长为 \(a\)，则 \(6a^2=24\)，得 \(a=2\). 内切球直径等于正方体棱长，所以半径 \(r=1\). 因而球的体积为 \(\dfrac43\pi r^3=\dfrac43\pi\mathrm{~cm}^3\)，选 A.
\end{solution}

\begin{problem}
“平面 \(\alpha\) 内有无穷多条直线都和直线 \(l\) 平行”是“\(l \parallel \alpha\)”的（\quad）

\choices
    {充分而不必要的条件}
    {必要而不充分的条件}
    {必要且充分的条件}
    {既不充分又不必要的条件}
\end{problem}

\begin{answer}
B
\end{answer}

\begin{solution}
若 \(l\parallel\alpha\)，则 \(l\) 的方向平行于平面 \(\alpha\)，平面内可作无穷多条与 \(l\) 平行的直线，所以题述条件是必要条件. 反之，平面内有无穷多条与 \(l\) 平行的直线只能说明 \(l\) 的方向平行于 \(\alpha\)，并不能排除 \(l\) 在 \(\alpha\) 内的情形，因此不充分，选 B.
\end{solution}

\begin{problem}
下列函数中，在区间 \((0,1)\) 上是增函数的只有（\quad）

\choices
    {\(y = -\sqrt{x}\)}
    {\(y = \frac{1}{x^3}\)}
    {\(y = \log_{\frac{1}{2}} x\)}
    {\(y = -x^2 + 2x + 1\)}
\end{problem}

\begin{answer}
D
\end{answer}

\begin{solution}
在 \((0,1)\) 上，\(-\sqrt{x}\) 与 \(x^{-3}\) 都递减，底数 \(\dfrac12\in(0,1)\) 的对数函数 \(\log_{1/2}x\) 也递减. 对 \(y=-x^2+2x+1\) 求导得 \(y'=2(1-x)>0\)，所以只有该函数递增，选 D.
\end{solution}

\begin{problem}
下列不等式中，与不等式 \(\frac{x - 3}{2 - x} \geqslant 0\) 同解的是（\quad）

\choices
    {\((x - 3)(2 - x) \geqslant 0\)}
    {\(\lg(x - 2) \leqslant 0\)}
    {\(\frac{2 - x}{x - 3} \geqslant 0\)}
    {\((x - 3)(2 - x) > 0\)}
\end{problem}

\begin{answer}
B
\end{answer}

\begin{solution}
原不等式的定义域为 \(x\ne2\). 在三个区间上分别判断符号：当 \(x<2\) 时，分子为负、分母为正，分式为负；当 \(2<x<3\) 时，分子、分母均为负，分式为正；当 \(x>3\) 时，分子为正、分母为负，分式为负；在 \(x=3\) 时分式为零. 所以原不等式的解集为 \(2<x\leqslant3\). 另一方面，\(\lg(x-2)\leqslant0\) 的定义域要求 \(x-2>0\)，且因常用对数函数递增，\(\lg(x-2)\leqslant0\) 等价于 \(0<x-2\leqslant1\)，解集同为 \(2<x\leqslant3\). 因此选 B.
\end{solution}

\begin{problem}
用 \(0\)，\(1\)，\(2\)，\(3\)，\(4\) 这五个数字组成没有重复数字的四位数，那么，在这些四位数中，是偶数的总共有（\quad）

\choices
    {\(120\) 个}
    {\(96\) 个}
    {\(60\) 个}
    {\(36\) 个}
\end{problem}

\begin{answer}
C
\end{answer}

\begin{solution}
末位为 \(0\) 时，首位有 \(4\) 种选法，其余两位有 \(\mathrm A_3^2\) 种，得 \(4\mathrm A_3^2=24\) 个. 末位为 \(2\) 或 \(4\) 时有 \(2\) 种选法，首位有 \(3\) 种，第二、三位依次有 \(3\)、\(2\) 种，得 \(2\times3\times3\times2=36\) 个. 总数为 \(24+36=60\)，故选 C.
\end{solution}

\begin{problem}
如果 \(\tan(\alpha + \beta) = \frac{2}{5}\)，\(\tan\left(\beta - \frac{\pi}{4}\right) = \frac{1}{4}\)，那么，\(\tan\left(\alpha + \frac{\pi}{4}\right)=\)（\quad）

\choices
    {\(\frac{13}{18}\)}
    {\(\frac{13}{22}\)}
    {\(\frac{3}{22}\)}
    {\(\frac{3}{18}\)}
\end{problem}

\begin{answer}
C
\end{answer}

\begin{solution}
由 \(\alpha+\dfrac\pi4=(\alpha+\beta)-\left(\beta-\dfrac\pi4\right)\)，设 \(A=\alpha+\beta\)、\(B=\beta-\dfrac\pi4\). 则
\(\tan(A-B)=\dfrac{\frac25-\frac14}{1+\frac25\cdot\frac14}=\dfrac{3/20}{22/20}=\dfrac3{22}\)，故选 C.
\end{solution}

\begin{problem}
函数值 \(\tan 224^{\circ}\)，\(\sin 136^{\circ}\)，\(\cos 310^{\circ}\) 的大小关系是（\quad）

\choices
    {\(\tan 224^{\circ} < \sin 136^{\circ} < \cos 310^{\circ}\)}
    {\(\sin 136^{\circ} < \cos 310^{\circ} < \tan 224^{\circ}\)}
    {\(\cos 310^{\circ} < \tan 224^{\circ} < \sin 136^{\circ}\)}
    {\(\cos 310^{\circ} < \sin 136^{\circ} < \tan 224^{\circ}\)}
\end{problem}

\begin{answer}
D
\end{answer}

\begin{solution}
有 \(\tan224^{\circ}=\tan44^{\circ}\)，且 \(\tan44^{\circ}>\sin44^{\circ}\). 又 \(\sin136^{\circ}=\sin44^{\circ}\)，\(\cos310^{\circ}=\cos50^{\circ}=\sin40^{\circ}\)，而正弦函数在 \([0,90^{\circ}]\) 上递增，所以 \(\cos310^{\circ}<\sin136^{\circ}<\tan224^{\circ}\). 故选 D.
\end{solution}

\begin{problem}
如果椭圆的两个焦点将长轴分成三等分，那么这个椭圆的两条准线间的距离是焦距的（\quad）

\choices
    {\(18\) 倍}
    {\(12\) 倍}
    {\(9\) 倍}
    {\(4\) 倍}
\end{problem}

\begin{answer}
C
\end{answer}

\begin{solution}
设椭圆长半轴为 \(a\)，焦半距为 \(c\). 两焦点把长轴分成三等份，因此 \(a-c=2c\)，得 \(a=3c\)，离心率 \(e=\dfrac ca=\dfrac13\). 两准线间距离为 \(\dfrac{2a}{e}\)，焦距为 \(2c\)，故其比值为 \(\dfrac{a}{ec}=\dfrac{a^2}{c^2}=9\)，选 C.
\end{solution}

\begin{problem}
记函数 \(y = \frac{1}{x}\) 的图象为 \(l_1\)，\(y = \arctan x\) 的图象为 \(l_2\)，那么 \(l_1\) 与 \(l_2\) 的交点个数是（\quad）

\choices
    {无穷多个}
    {\(2\) 个}
    {\(1\) 个}
    {\(0\) 个}
\end{problem}

\begin{answer}
B
\end{answer}

\begin{solution}
交点满足 \(\dfrac1x=\arctan x\)，其中 \(x\ne0\)，等价于 \(g(x)=x\arctan x=1\). 在 \(x>0\) 上，\(g'(x)=\arctan x+\dfrac{x}{1+x^2}>0\)，所以 \(g\) 严格递增；并且 \(\lim_{x\to0^+}g(x)=0\)，当 \(x\geqslant1\) 时 \(g(x)\geqslant\dfrac\pi4x\)，故 \(\lim_{x\to+\infty}g(x)=+\infty\). 由连续性，\(g(x)=1\) 在正半轴上恰有一个根. 又 \(g(-x)=g(x)\)，所以有且仅有一个对应的负根，交点共有 \(2\) 个，选 B.
\end{solution}

\begin{problem}
如果实数 \(x\)，\(y\) 满足 \(x^2 + y^2 = 1\)，那么 \((1 - xy)(1 + xy)\) 有（\quad）

\choices
    {最小值 \(\frac{1}{2}\) 和最大值 \(1\)}
    {最大值 \(1\) 和最小值 \(\frac{3}{4}\)}
    {最小值 \(\frac{3}{4}\) 而无最大值}
    {最大值 \(1\) 而无最小值}
\end{problem}

\begin{answer}
B
\end{answer}

\begin{solution}
由 \(x^2+y^2=1\) 及基本不等式，\(|xy|\leqslant\dfrac{x^2+y^2}{2}=\dfrac12\)，所以 \(0\leqslant x^2y^2\leqslant\dfrac14\). 因而
\((1-xy)(1+xy)=1-x^2y^2\) 的最大值为 \(1\)，最小值为 \(\dfrac34\)，且分别在 \(xy=0\)、\(|x|=|y|=\dfrac1{\sqrt2}\) 时取得，选 B.
\end{solution}

\begin{problem}
曲线 \(y = 1 + \sqrt{4 - x^2}\)（\(-2 \leqslant x \leqslant 2\)）与直线 \(y = k(x - 2) + 4\) 有两个交点时，实数 \(k\) 的取值范围是（\quad）

\choices
    {\(\left(\frac{5}{12}, \frac{3}{4}\right]\)}
    {\(\left(\frac{5}{12}, +\infty\right)\)}
    {\(\left(\frac{1}{3}, \frac{3}{4}\right]\)}
    {\(\left(0, \frac{5}{12}\right)\)}
\end{problem}

\begin{answer}
A
\end{answer}

\begin{solution}
圆弧满足 \(x^2+(y-1)^2=4\) 且 \(y\geqslant1\). 代入直线 \(y=k(x-2)+4\)，得
\( (1+k^2)x^2+2k(3-2k)x+(3-2k)^2-4=0\). 其判别式为 \(4(12k-5)\)，所以直线与整圆有两个交点须有 \(k>\dfrac5{12}\)，这时 \(k>0\). 设两个交点的横坐标为 \(x_1<x_2\). 在线段所在直线上，\(y\geqslant1\) 等价于 \(x\geqslant x_0=2-\dfrac3k\).

当 \(\dfrac5{12}<k\leqslant\dfrac34\) 时，\(x_0\leqslant-2\)，而圆上横坐标都在 \([-2,2]\) 内，所以两个交点都在上半圆弧上. 当 \(k>\dfrac34\) 时，\(-2<x_0<2\)，且把 \(x=x_0\) 代入圆的左端表达式得 \(x_0^2-4<0\)，故 \(x_0\) 位于两个根之间，即 \(x_1<x_0<x_2\)，左侧交点在下半圆，只有右侧交点在上半圆. 因而所求范围为 \(\left(\dfrac5{12},\dfrac34\right]\)，选 A.
\end{solution}

\begin{problem}
设 \(f(x)\) 是 \(\R\) 上的奇函数，且当 \(x \in [0,+\infty)\) 时，\(f(x) = x\left(1 + \sqrt[3]{x}\right)\)，那么 \(x \in (-\infty,0)\) 时，\(f(x)=\)（\quad）

\choices
    {\(-x\left(1 + \sqrt[3]{x}\right)\)}
    {\(x\left(1 + \sqrt[3]{x}\right)\)}
    {\(-x\left(1 - \sqrt[3]{x}\right)\)}
    {\(x\left(1 - \sqrt[3]{x}\right)\)}
\end{problem}

\begin{answer}
D
\end{answer}

\begin{solution}
当 \(x<0\) 时，\(-x>0\). 由奇函数性质，\(f(x)=-f(-x)=-\left[(-x)\left(1+\sqrt[3]{-x}\right)\right]\). 因为 \(\sqrt[3]{-x}=-\sqrt[3]{x}\)，所以 \(f(x)=x\left(1-\sqrt[3]{x}\right)\)，选 D.
\end{solution}

\begin{problem}
已知 \(y = f(x)\) 的图象如右图，那么 \(f(x)=\)（\quad）

\choices
    {\(\sqrt{x^2 - 2|x| + 1}\)}
    {\(x^2 - 2|x| + 1\)}
    {\(|x^2 - 1|\)}
    {\(\sqrt{x^2 - 2x + 1}\)}
\end{problem}

\begin{answer}
A
\end{answer}

\begin{solution}
原始试卷在本题题干中写有“如右图”，但当前扫描件在题干与选项之间未包含该条件图，因此无法从现有题面独立读取图象特征并完成唯一判定. 将四个选项代数化，依次得到 \(||x|-1|\)、\((|x|-1)^2\)、\(|x^2-1|\) 和 \(|x-1|\). 前三项都是偶函数，都非负，且都在 \(x=\pm1\) 处取值 \(0\)、在 \(x=0\) 处取值 \(1\)；它们只能由中间和外部曲线的具体形状区分. 第四项不是偶函数，但仅凭现有文字仍不能在前三项中确定选项. 原始试卷答案页给出第一项，但答案页不能补回缺失的图象条件，故本题保留阻塞，不猜测图形并修改答案.
\end{solution}

\begin{problem}
\(\sqrt{1 + \sin 10}=\)（\quad）

\choices
    {\(\cos 5 + \sin 5\)}
    {\(-\cos 5 - \sin 5\)}
    {\(\sqrt{2}\cos 5\)}
    {\(\cos 5 - \sin 5\)}
\end{problem}

\begin{answer}
B
\end{answer}

\begin{solution}
这里角的单位为弧度. 利用半角公式，\(1+\sin10=\sin^25+2\sin5\cos5+\cos^25=(\sin5+\cos5)^2\). 因为 \(\dfrac{3\pi}{2}<5<2\pi\)，且 \(\pi<5-\dfrac\pi4<\dfrac{3\pi}{2}\)，所以 \(\cos5+\sin5=\sqrt2\cos\left(5-\dfrac\pi4\right)<0\). 故平方根取其相反数，即 \(\sqrt{1+\sin10}=-\cos5-\sin5\)，选 B.
\end{solution}

\begin{problem}
一个三棱锥，如果它的底面是直角三角形，那么它的三个侧面（\quad）

\choices
    {必然都是非直角三角形}
    {至多只能有一个是直角三角形}
    {至多只能有两个是直角三角形}
    {可能都是直角三角形}
\end{problem}

\begin{answer}
D
\end{answer}

\begin{solution}
构造四面体 \(P-ABC\)：令底面三角形 \(ABC\) 在同一平面内且 \(\angle ABC=90^{\circ}\)，再取 \(PA\perp\) 平面 \(ABC\). 于是 \(PA\perp AB\) 且 \(PA\perp AC\)，所以侧面 \(PAB\)、\(PAC\) 都是直角三角形；又因为 \(AB\perp BC\)，且 \(PA\perp BC\)，由 \(\overrightarrow{PB}=\overrightarrow{PA}+\overrightarrow{AB}\) 得 \(PB\perp BC\)，因此侧面 \(PBC\) 也是直角三角形. 所以三个侧面可能都是直角三角形，选 D.
\end{solution}

\section{填空题}

\begin{problem}
把极坐标方程 \(\rho = 4\sin \theta\) 化成直角坐标方程得\fillinblank{}.
\end{problem}

\begin{answer}
\(x^2 + y^2 - 4y = 0\)
\end{answer}

\begin{solution}
由极坐标与直角坐标的关系 \(x=\rho\cos\theta\)、\(y=\rho\sin\theta\) 及 \(\rho^2=x^2+y^2\)，原方程两边乘以 \(\rho\) 得 \(\rho^2=4\rho\sin\theta=4y\). 因而直角坐标方程为 \(x^2+y^2-4y=0\).
\end{solution}

\begin{problem}
函数 \(y = |\cos x|\) 的最小正周期是\fillinblank{}.
\end{problem}

\begin{answer}
\(\pi\)
\end{answer}

\begin{solution}
因为 \(|\cos(x+\pi)|=|{-\cos x}|=|\cos x|\)，所以 \(\pi\) 是周期. 若 \(T>0\) 是任一周期，则由 \(|\cos(\frac\pi2+T)|=|\cos\frac\pi2|=0\)，可知 \(\frac\pi2+T=\frac\pi2+k\pi\)，其中 \(k\) 是整数；由 \(T>0\) 得 \(k\geqslant1\)，所以 \(T\geqslant\pi\). 因此最小正周期为 \(\pi\).
\end{solution}

\begin{problem}
复数 \(-\sqrt{6} + \sqrt{2}\i\) 的三角形式是\fillinblank{}.
\end{problem}

\begin{answer}
\(2\sqrt{2}\left(\cos \frac{5\pi}{6} + \i \sin \frac{5\pi}{6}\right)\)
\end{answer}

\begin{solution}
复数的模为 \(\sqrt{(-\sqrt6)^2+(\sqrt2)^2}=2\sqrt2\). 其对应点在第二象限，且 \(\cos\varphi=-\dfrac{\sqrt3}{2}\)、\(\sin\varphi=\dfrac12\)，所以 \(\varphi=\dfrac{5\pi}{6}\). 因而三角形式为 \(2\sqrt2\left(\cos\dfrac{5\pi}{6}+\i\sin\dfrac{5\pi}{6}\right)\).
\end{solution}

\begin{problem}
\(\displaystyle\lim_{n \to \infty}\frac{1 - 2 + 3 - 4 + \cdots + (2n - 1) - 2n}{(n - 1)\left(\frac{1}{n} + 1\right)}\) 的值是\fillinblank{}.
\end{problem}

\begin{answer}
\(-1\)
\end{answer}

\begin{solution}
分子按相邻两项分组为 \((1-2)+(3-4)+\cdots+((2n-1)-2n)=-n\). 分母为 \((n-1)\left(1+\dfrac1n\right)=\dfrac{n^2-1}{n}\)，所以原式等于 \(-\dfrac{n^2}{n^2-1}\)，其极限为 \(-1\).
\end{solution}

\begin{problem}
在 \((1 + x + x^2)(1 - x)^{10}\) 的展开式中，\(x^5\) 的系数是\fillinblank{}.
\end{problem}

\begin{answer}
\(-162\)
\end{answer}

\begin{solution}
取 \(x^5\) 时，三项分别贡献 \(\mathrm C_{10}^{5}(-1)^5\)、\(\mathrm C_{10}^{4}(-1)^4\)、\(\mathrm C_{10}^{3}(-1)^3\). 因而系数为 \(-252+210-120=-162\).
\end{solution}

\section{解答题}

\begin{problem}
如图，在四棱锥 \(P-ABCD\) 中，侧面 \(PCD\) 是边长为 \(2\mathrm{~cm}\) 的等边三角形，且与底面垂直，而底面 \(ABCD\) 是面积为 \(2\sqrt{3}\mathrm{~cm}^2\) 的菱形，\(\angle ADC\) 是锐角.
\begin{enumerate}
\item 求这个四棱锥的体积；
\item 求证：\(PA \perp CD\).
\end{enumerate}

\bitmapfigure[max width=0.42\linewidth]{img/1988/guangdong_science/q26_fig1.png}
\end{problem}

\begin{solution}
\begin{enumerate}
\item 过 \(P\) 作 \(PE \perp CD\)，垂足是 \(E\). 因为 \(CD\) 是侧面 \(PCD\) 与底面 \(ABCD\) 的交线，侧面 \(PCD \perp\) 底面 \(ABCD\)，且 \(PE\) 在侧面 \(PCD\) 内，所以 \(PE \perp\) 底面 \(ABCD\). 又因为 \(\triangle PCD\) 是边长为 \(2\mathrm{~cm}\) 的等边三角形，且 \(PE \perp CD\)，所以 \(PE=\sqrt{3}\mathrm{~cm}\). 依设 \(S_{\text{底面}}=2\sqrt{3}\mathrm{~cm}^2\)，因此四棱锥的体积为

\[
V=\frac{1}{3}PE\cdot S_{\text{底面}}=\frac{1}{3}\times\sqrt{3}\times2\sqrt{3}=2\mathrm{~cm}^3
\]

\item 记 \(\angle ADC=\theta\). 因为底面 \(ABCD\) 是菱形，所以 \(S_{\text{底面}}=CD^2\sin\theta\). 由题设 \(CD=2\mathrm{~cm}\)，\(S_{\text{底面}}=2\sqrt{3}\mathrm{~cm}^2\)，得 \(\sin\theta=\frac{2\sqrt{3}}{4}=\frac{\sqrt{3}}{2}\). 又因为 \(\theta\) 是锐角，所以 \(\theta=60^{\circ}\). 连结 \(AC\)，得 \(\triangle ADC\) 是等边三角形. 又因为 \(\triangle PCD\) 也是等边三角形，且 \(PE\perp CD\)，所以 \(E\) 是 \(CD\) 的中点. 连结 \(AE\)，则 \(AE\perp CD\). 又 \(AE\cap PE=E\)，故 \(CD\perp\) 平面 \(PAE\). 因为 \(PA\) 在平面 \(PAE\) 内，所以 \(PA\perp CD\).
\end{enumerate}
\end{solution}

\begin{problem}
已知双曲线的方程是 \(16x^2 - 9y^2 = 144\).
\begin{enumerate}
\item 求这双曲线的焦点坐标，离心率和渐近线方程；
\item 设 \(F_1\) 和 \(F_2\) 是这双曲线的左，右焦点，点 \(P\) 在这双曲线上，且 \(\left|PF_1\right|\left|PF_2\right| = 32\)，求 \(\angle F_1PF_2\) 的大小.
\end{enumerate}
\end{problem}

\begin{solution}
\begin{enumerate}
\item 把双曲线方程化为标准方程式 \(\frac{x^2}{3^2}-\frac{y^2}{4^2}=1\). 因此实半轴长 \(a=3\)，虚半轴长 \(b=4\)，中心在坐标原点，实轴在 \(x\) 轴上. 由 \(c=\sqrt{a^2+b^2}=\sqrt{9+16}=5\)，得焦点坐标为 \(F_1(-5,0)\) 和 \(F_2(5,0)\)，离心率 \(e=\frac{c}{a}=\frac{5}{3}\)，渐近线方程为 \(y=\pm\frac{b}{a}x=\pm\frac{4}{3}x\).
\end{enumerate}

\solutionfigure{\bitmapfigure[max width=0.54\linewidth]{img/1988/guangdong_science/q27_solution_fig1.png}}

\begin{enumerate}
\setcounter{enumi}{1}
\item 设 \(\left|PF_1\right|=d_1\)，\(\left|PF_2\right|=d_2\). 由题设 \(d_1d_2=32\). 又由双曲线的几何性质知 \(\left|d_1-d_2\right|=2a=6\)，所以 \(d_1^2+d_2^2-2d_1d_2=36\)，即 \(d_1^2+d_2^2=36+2d_1d_2=100\). 又 \(\left|F_1F_2\right|=2c=10\). 因此
\[
\left|F_1F_2\right|^2=100=d_1^2+d_2^2=\left|PF_1\right|^2+\left|PF_2\right|^2
\]
于是 \(\triangle PF_1F_2\) 是以 \(F_1F_2\) 为斜边的直角三角形，所以 \(\angle F_1PF_2=90^{\circ}\).
\end{enumerate}
\end{solution}

\begin{problem}
设 \(0 < \theta < 2\pi\)，复数 \(z = 1 - \cos \theta + \i \sin \theta\)，\(u = a^2 + a\i\)，且 \(zu\) 是纯虚数. 这里，\(\i\) 是虚数单位，\(a\) 是实数.
\begin{enumerate}
\item 求复数 \(u\) 的辐角主值 \(\arg u\)（用 \(\theta\) 的代数式表示）；
\item 记 \(\omega = z^2 + u^2 + 2zu\)，试问：\(\omega\) 可能是正实数吗？为什么？
\end{enumerate}
\end{problem}

\begin{solution}
\begin{enumerate}
\item 因为 \(zu=\left[a^2(1-\cos\theta)-a\sin\theta\right]+\left[a^2\sin\theta+a(1-\cos\theta)\right]\i\)，而 \(zu\) 是纯虚数，所以 \(a^2\sin\theta+a(1-\cos\theta)\neq0\)，\(a^2(1-\cos\theta)-a\sin\theta=0\). 由前一式知 \(a\neq0\). 又因为 \(0<\theta<2\pi\)，所以 \(1-\cos\theta\neq0\). 因此
\[
a=\frac{\sin\theta}{1-\cos\theta}=\frac{2\cos\frac{\theta}{2}\sin\frac{\theta}{2}}{2\sin^2\frac{\theta}{2}}=\cot\frac{\theta}{2}
\]
又由上式知 \(\theta\neq\pi\). 所以
\[
\tan(\arg u)=\frac{a}{a^2}=\frac{1}{a}=\tan\frac{\theta}{2}
\]
因为 \(0\leqslant\arg u<2\pi\)，所以当 \(0<\theta<\pi\) 时，复平面上与复数 \(u=\cot^2\frac{\theta}{2}+\cot\frac{\theta}{2}\i\) 对应的点在第一象限，得 \(\arg u=\frac{\theta}{2}\). 当 \(\pi<\theta<2\pi\) 时，该点在第四象限，得 \(\arg u=\pi+\frac{\theta}{2}\).

\item \(\omega\) 不可能是正实数. 因为 \(\omega=(z+u)^2\)，且
\[
z+u=1-\cos\theta+a^2+(a+\sin\theta)\i
\]
记 \(z+u=R+I\i\)，则 \(R=1-\cos\theta+a^2>0\). 若 \(\omega=(z+u)^2\) 是正实数，则其虚部 \(2RI\) 必须为零；由 \(R>0\) 得 \(I=0\)，即 \(a+\sin\theta=0\)，所以 \(a=-\sin\theta\). 但由上题已得 \(a=\frac{\sin\theta}{1-\cos\theta}\)，且 \(\sin\theta\neq0\). 因此 \(-\sin\theta=\frac{\sin\theta}{1-\cos\theta}\)，所以 \(-1=\frac{1}{1-\cos\theta}\)，解得 \(\cos\theta=2\)，这不可能. 因此 \(\omega\) 不能是正实数.
\end{enumerate}
\end{solution}

\begin{problem}
设 \(a > 0\)，\(a \neq 1\)，\(f(x) = \log_a\left(x + \sqrt{x^2 - 1}\right)\)（\(x \geqslant 1\)）.
\begin{enumerate}
\item 求函数 \(f(x)\) 的反函数 \(f^{-1}(x)\) 和这个反函数的定义域；
\item 记 \(b_n = f^{-1}(n)\)，\(S_n = b_1 + b_2 + \cdots + b_n\)，求证：当 \(\frac{1}{2} < a < 2\) 时，对任意自然数 \(n\)，都有 \(S_n < 2^n - \left(\frac{\sqrt{2}}{2}\right)^n\).
\end{enumerate}
\end{problem}

\begin{solution}
\begin{enumerate}
\item 由 \(y=\log_a\left(x+\sqrt{x^2-1}\right)\)（\(x\geqslant1\)）得 \(x+\sqrt{x^2-1}=a^y\). 又因为
\(x+\sqrt{x^2-1}=\frac{x^2-(x^2-1)}{x-\sqrt{x^2-1}}=\frac{1}{x-\sqrt{x^2-1}}\)，所以 \(x-\sqrt{x^2-1}=a^{-y}\). 由这两式可得 \(x=\frac{a^y+a^{-y}}{2}\)，因此反函数为 \(f^{-1}(x)=\frac{a^x+a^{-x}}{2}\).

令 \(h(x)=x+\sqrt{x^2-1}\). 在 \([1,+\infty)\) 上，\(h(1)=1\)，且 \(h\) 严格递增并趋于 \(+\infty\)，所以 \(h\) 的值域是 \([1,+\infty)\). 因而函数 \(f\) 的值域与 \(a\) 的取值有关：当 \(a>1\) 时，\(f\) 的值域是 \(\left[0,+\infty\right)\)；当 \(0<a<1\) 时，\(f\) 的值域是 \(\left(-\infty,0\right]\). 反函数的定义域分别就是这两个区间.

\item 按题面给出的 \(\frac{1}{2}<a<2\)，若 \(\frac{1}{2}<a<1\)，则上问所得反函数定义域为 \((-\infty,0]\)，所以 \(b_1=f^{-1}(1)\) 无定义，第二问在这段参数范围内并没有定义. 因而以下证明针对使所有 \(b_n\) 有定义的参数范围 \(1<a<2\)；这也明确指出了原题参数范围的题面矛盾. 在此范围内，\(1<(2a)^n\)，且 \(a^n<2^n\)，因而
\[
\begin{gathered}
\frac{1}{2}\left[(2^n+2^{-n})-(a^n+a^{-n})\right]=\frac{1}{2(2a)^n}\left[2^{2n}a^n+a^n-a^{2n}2^n-2^n\right] \\
=\frac{\left[(2a)^n-1\right]\left(2^n-a^n\right)}{2(2a)^n}>0.
\end{gathered}
\]
所以 \(\frac{1}{2}(a^n+a^{-n})<\frac{1}{2}(2^n+2^{-n})\). 由上题知 \(b_n=f^{-1}(n)=\frac{1}{2}(a^n+a^{-n})\)，因此
\[
\begin{gathered}
S_n<\frac{1}{2}(2+2^{-1})+\frac{1}{2}(2^2+2^{-2})+\cdots+\frac{1}{2}(2^n+2^{-n}) \\
=\frac{1}{2}(2+2^2+\cdots+2^n)+\frac{1}{2}(2^{-1}+2^{-2}+\cdots+2^{-n}) \\
=2^n-1+\frac{1}{4}\left[\frac{1-2^{-n}}{1-2^{-1}}\right]=2^n-\frac{1}{2}(1+2^{-n}).
\end{gathered}
\]
对任意自然数 \(n\geqslant1\)，都有 \(2^{-n}\neq1\)，所以 \(\frac{1}{2}(1+2^{-n})>\sqrt{2^{-n}}=\left(\frac{\sqrt{2}}{2}\right)^n\)，于是 \(2^n-\frac{1}{2}(1+2^{-n})<2^n-\left(\frac{\sqrt{2}}{2}\right)^n\). 故对任意自然数 \(n\geqslant1\)，都有 \(S_n<2^n-\left(\frac{\sqrt{2}}{2}\right)^n\)，即证.
\end{enumerate}
\end{solution}
